\begin{lemma}[Terminal PE tree interchangeability away from delta]
Let $T$ be a positive-rung extended Tashkinov tree constructed from a
$k$-triple $G,e_0,\varphi$, with arbitrary certified input coloring,
terminal ladder tree $T_n$, terminal generating coloring $\varphi_n$,
and terminal stored-color set $D_n$.  Assume that $T$ has maximum
property and that its terminal extension is a PE.  Let $R_n$ and
$T_n\vee R_n$ be the closure and exact ordered union in a hierarchy
base under $\varphi_n$ (respectively `base.closure' and `base.core'
in Lean).

Assume, uniformly over every certified input coloring and generating
coloring, that every positive-rung ETT $T'$ of rung number $s<n$ with
maximum property satisfies all six lower-rung assertions:
\begin{enumerate}
\item $V(T')$ is elementary under its generating coloring;
\item under every proper coloring leaving exactly $e_0$ uncolored and
stable on its final ladder closure, that closure is interchangeable for
two colors whenever at least one of them is missing there;
\item stability descends at each PE support to the prefix before the
support and to its first-appearance ladder stage, and stored pairs at
distinct PE slots with distinct supports and the same first-appearance
index are disjoint;
\item at a terminal PE, the stored two-color component through the
support meets the terminal ladder tree only at that support;
\item every nonminimum endpoint of a defective boundary color precedes
each later PE support or SE extension vertex allowed by its
first-appearance index; and
\item every proper terminal stable coloring leaving exactly $e_0$
uncolored is a mod-coloring of the ETT.
\end{enumerate}

Let $\sigma_n$ be proper, leave exactly $e_0$ uncolored, and be
$T_n\oplus R_n,D_n,\varphi_n$-stable.

Let $v_n\in V(G)$ and $\gamma_n,\delta_n\in[k]$ be exposed by the
exact terminal option-slot equalities
\[
 \mathcal T.\mathrm{chosenVertex}(n)=\mathsf{some}(v_n),\qquad
 \mathcal T.\mathrm{gamma}(n)=\mathsf{some}(\gamma_n),\qquad
 \mathcal T.\mathrm{delta}(n)=\mathsf{some}(\delta_n).
\]
For arbitrary $\alpha,\beta\in[k]$, if
$\alpha\in\overline\sigma_n(T_n)$,
$\alpha\ne\delta_n$, and $\beta\ne\delta_n$, then
$\alpha,\beta$ are $T_n$-interchangeable under $\sigma_n$.
All boundaries, trees, and two-color components use the original edge
type, so distinct parallel edges retain their identities.
\end{lemma}
\begin{proof}
Write $X=V(T_n)$, $Y=V(R_n)$, and
$U=V(T_n\vee R_n)$, with vertex sets understood as in
\noderef{TreeSequenceVertexSet}.  The generating-color field of
\noderef{ExtendedTashkinovTree} identifies $\varphi_n$ with the last
coloring of the certified series.  Eliminate the last certificate in
\noderef{TashkinovSeries} using the terminal-PE hypothesis.  The PE
constructor of \noderef{TashkinovSeriesStep}, together with the three
assumed last-slot equalities, identifies its support and colors with the
displayed $v_n,\gamma_n,\delta_n$ and supplies the exact connecting-edge
identity $f_n$.  It also supplies properness and exact
one-uncoloredness of $\varphi_n$, and the equations recording the
terminal $\gamma_n,\delta_n$-switch.  No endpoint pair is substituted
for $f_n$.

The parallel clause of \noderef{HierarchyBase} says that $R_n$ is a
$\varphi_n$-closure of the ordered prefix $T_n(v_n)$, and supplies
an \noderef{OrderedTreeUnion} certificate with
\[
                         U=X\cup Y.                       \tag{1}
\]
It also says that the missing-color set of $T_n$ under the base
coloring is its missing-color set under $\varphi_n$.
The two conjuncts of \noderef{CombinedStableColoring}, with stability
expanded as in \noderef{StableColoring}, say that
$\sigma_n$ is $(T_n,D_n,\varphi_n)$-stable and simultaneously
$(R_n,\varnothing,\varphi_n)$-stable.  Applying
\noderef{CombinedStableClosurePreserved} gives, for every displayed edge
identity $g$ of either tree,
\[
             \sigma_n(g)=\varphi_n(g),                    \tag{2}
\]
and proves that the same ordered tree $R_n$, with the same edge
identities, is a $\sigma_n$-closure of $T_n(v_n)$.

We record three preliminary facts, with their hypotheses checked at the
coloring at which they will be used.

First, $U$ is elementary under $\sigma_n$.  Indeed, retain the same
core, closure, and support position from the hierarchy base.  Equality
of missing-color sets on $T_n$ is the vertex clause of the first stable
relation.  The preceding paragraph supplies the required
$\sigma_n$-closure, and the old ordered-union certificate is unchanged.
These data form a hierarchy base under $\sigma_n$.  The first stable
relation, properness, exact one-uncoloredness, positivity, maximum
property, the terminal-PE equation, and the displayed six lower-rung
assertions are precisely the hypotheses of
\noderef{ParallelClosureUnionElementary}.  Hence
\[
                  U\text{ is elementary under }\sigma_n.   \tag{3}
\]

The vertex clauses of the two stable relations identify both constituent
missing-color sets under $\varphi_n$ and $\sigma_n$.  The PE constructor
of \noderef{TashkinovSeriesStep} says that $\varphi_n$ is proper and
leaves exactly $e_0$ uncolored; the same two facts for $\sigma_n$ are
hypotheses.  We may therefore apply the symmetry clause of \noderef{sc2}
separately to the two constituent stable relations.  It gives the reverse
relations saying that $\varphi_n$ is $(T_n,D_n,\sigma_n)$-stable and is
also $(R_n,\varnothing,\sigma_n)$-stable.

Third, we establish the following restricted terminal-tree assertion:
\[
 \begin{split}
 a\in\overline\sigma_n(T_n),\quad
 a,b\notin\{\gamma_n,\delta_n\}
 \quad\Longrightarrow\quad
 a,b\text{ are }T_n\text{-interchangeable under }\sigma_n.
                                                               \tag{6}
 \end{split}
\]
Apply \noderef{ExtensionBase} with the six lower-rung assertions, the
given maximum property, properness, exact one-uncoloredness, and the
first stable relation.  Its terminal-PE conclusion says that the
$\gamma_n,\delta_n$-component through $v_n$ meets $X$ exactly at
$v_n$.  Thus all hypotheses of
\noderef{ParallelExtensionUndoStable} hold.  If
\[
 \tau=\sigma_n/P_{v_n}(\gamma_n,\delta_n,\sigma_n),
\]
the PE certificate gives $\gamma_n\ne\delta_n$: the rule
\noderef{ParallelExtensionColor} makes $\gamma_n$ missing at $v_n$ under
the selected coloring, while \noderef{MaximumDefectiveChoice} supplies the
exact edge identity $f_n$ incident with $v_n$ and colored $\delta_n$.
Thus \noderef{PresentColors} and \noderef{MissingColors} put the two colors
on opposite sides of the same complement.  Hence
\noderef{TwoColorComponentSwapProper},
applied to the stated properness of $\sigma_n$, proves that $\tau$ is
proper.  To verify its uncolored edge exactly, unfold
\noderef{IsOneEdgeUncolored}.  The value $\sigma_n(e_0)=\mathsf{none}$
is fixed by the definition \noderef{TwoColorComponentSwap}, so
$\tau(e_0)=\mathsf{none}$.  If $g\ne e_0$, the same hypothesis supplies
a color $d$ with $\sigma_n(g)=\mathsf{some}(d)$.  The swap either fixes
this value or, when the edge-node of $g$ is in the selected component and
$d\in\{\gamma_n,\delta_n\}$, replaces it by the other displayed color.
In every case $\tau(g)$ is still a colored value.  Thus $\tau$ leaves
exactly $e_0$ uncolored.

The actual conclusion of \noderef{ParallelExtensionUndoStable} now makes
$\tau$ stable on $T_n$ relative to the preceding stored set and preceding
generating coloring.  Since $(a,b)$ avoid the two switched colors,
\noderef{TwoColorComponentSwap} shows edge by edge that the
$(a,b)$-incidence graph, its components, and the missingness of $a$ on
$T_n$ are identical for $\tau$ and $\sigma_n$.
When $n>1$, \noderef{CertifiedETTPredecessor} supplies the actual
rung-$n-1$ ETT with inherited maximum property, terminal tree exactly
$T_n$, preceding generating coloring, and all retained certificate
arrays equal to the corresponding arrays of $T$.  Lower assertion (ii)
applied to that predecessor and to $\tau$ gives interchangeability on
its final ladder closure $T_n$.  When $n=1$, the same conclusion is the
interchangeability clause of \noderef{TechnicalTheoremBase}.  Transport across the identical
incidence graphs proves (6).

Before proving (iii), we strengthen (6) to the paper's observation
\[
 \begin{split}
 a\in\overline\sigma_n(T_n),\qquad
 a\ne\delta_n,\quad b\ne\delta_n
 \quad\Longrightarrow\quad
 a,b\text{ are }T_n\text{-interchangeable}.                 \tag{8}
 \end{split}
\]
The unconditional clause of \noderef{HkuParallelBoundaryStructure} gives
\[
 \partial_{\varphi_n,\gamma_n}(T_n)=\{f_n\}.                  \tag{8a}
\]
For the conditional clause of the same lemma, pre-step elementarity is
available from the identical predecessor argument just used: lower
assertion (i) applies to the certified predecessor when $n>1$, and
\noderef{TechnicalTheoremBase} applies when $n=1$.  Since $\gamma_n\ne\delta_n$, if $\gamma_n$ were
missing on $T_n$ under $\varphi_n$, that clause would make the left side
of (8a) empty, contradicting the membership of the exact identity $f_n$
in its singleton right side.  Therefore
$\gamma_n\notin\overline\varphi_n(T_n)$.  The PE constructor of
\noderef{TashkinovSeriesStep} gives the exact current stored-pair equation
$S_n=\{\delta_n,\gamma_n\}$.  In the set difference defining $D_n$ in
\noderef{TashkinovStoredColors}, this puts $\gamma_n$ in the cumulative
stored-color union, and the just-proved nonmissingness prevents its
removal.  Hence $\gamma_n\in D_n$.

Use now both directions of the $T_n$-stable relation: the forward
relation is a constituent of \noderef{CombinedStableColoring}, and its
reverse was obtained from \noderef{sc2} above.  Since $\gamma_n\in D_n$,
their edge clauses apply whether a tree-incident edge is
$\gamma_n$-colored under $\varphi_n$ or under $\sigma_n$.  Exact boundary
membership is coloring-independent, so the two inclusions obtained from
the forward and reverse relations transport (8a) to
\[
 \partial_{\sigma_n,\gamma_n}(T_n)=\{f_n\}.                  \tag{8b}
\]

Now take $a$ as in (8).  The stable vertex clause makes it missing on
$T_n$ under $\varphi_n$.  Since $a\ne\delta_n$, the conditional clause of
\noderef{HkuParallelBoundaryStructure} makes its $\varphi_n$-boundary
empty.  If its $\sigma_n$-boundary contained an exact identity, the edge
clause of the reverse $T_n$-stable relation would give that identity color
$a$ under $\varphi_n$, a contradiction.  Thus the $a$-boundary under
$\sigma_n$ is empty.  Equation (8b) then gives $a\ne\gamma_n$.  If
$b\ne\gamma_n$, (8) is (6).  If $b=\gamma_n$, then (3), unfolded using
\noderef{ColoredVertexSetElementary}, makes $X\subseteq U$ elementary, the
$a$-boundary of $X$ is empty, and the $\gamma_n$-boundary contains
only $f_n$.  Properness and
\noderef{TwoColorComponentSpanningPath} then allow at most one paper
$(a,\gamma_n)$-path component to meet $X$.  For two nontrivial components,
the spanning paths would require two distinct boundary identities or give
a repeated missing color in $X$.  If one component were a singleton at a
tree vertex, both colors would be missing there; any second singleton or
any tree end of a trapped nontrivial component would repeat a missing
color, while a nontrivial component crossing $X$ cannot have both ends
outside through the sole boundary identity $f_n$.  Thus a singleton cannot
coexist with a second path component meeting $X$.  This proves (8).
\end{proof}
